% Figure 1. Continuation connectivity and invariant projection.
\begin{figure*}[t]
\centering
\begin{tikzpicture}[x=1cm,y=1cm,line cap=round,line join=round,
  >=Stealth,every node/.style={font=\footnotesize}]
  % ---- (a) schematic invariant projection ----
  \begin{scope}
    \node[font=\small,anchor=north west] at (-0.10,4.35) {(a)};
    \draw[line width=0.8pt,->] (0,0) -- (7.30,0)
      node[below left,font=\small] {$T_{\rm si}$};
    \draw[line width=0.8pt,->] (0,0) -- (0,4.15)
      node[above left,font=\small] {$L_{\rm si}$};
    \node[font=\scriptsize,anchor=north east] at (7.05,4.02)
      {schematic};

    \fill[black!12]
      (0.45,3.35) .. controls (2.25,3.65) and (4.45,3.10) .. (6.55,2.18)
      -- (6.55,1.94)
      .. controls (4.45,2.86) and (2.25,3.41) .. (0.45,3.11) -- cycle;
    \draw[line width=1.05pt,black!70]
      (0.45,3.23) .. controls (2.25,3.53) and (4.45,2.98) .. (6.55,2.06);

    \fill[black!28]
      (0.45,1.05) .. controls (2.20,0.78) and (4.35,1.26) .. (6.55,1.88)
      -- (6.55,2.12)
      .. controls (4.35,1.50) and (2.20,1.02) .. (0.45,1.29) -- cycle;
    \draw[line width=1.05pt,black]
      (0.45,1.17) .. controls (2.20,0.90) and (4.35,1.38) .. (6.55,2.00);

    % A and B represent the documented near-coincident invariant pair.
    \fill[black] (5.83,2.29) circle (2.7pt);
    \node[anchor=south east] at (5.73,2.47) {A};
    \draw[line width=0.9pt,black,fill=white]
      (5.88,1.88) rectangle ++(0.28,0.28);
    \node[anchor=north west] at (6.20,1.91) {B};
    \draw[line width=0.65pt,black!55,dashed] (5.83,2.29) -- (6.02,2.02);

    \node[font=\scriptsize,anchor=west] at (0.55,3.72) {projected branch};
    \node[font=\scriptsize,anchor=west] at (0.55,0.62) {projected branch};
  \end{scope}

  % ---- (b) straight mass segment used in the documented continuation ----
  \begin{scope}[shift={(8.35,0.05)}]
    \node[font=\small,anchor=north west] at (-0.15,4.30) {(b)};
    % window m1 in [1.05,1.115], m2 in [1.04,1.145]
    \draw[line width=0.8pt] (0,0) rectangle (6.30,3.85);
    \foreach \x/\lab in {0.000/1.05,2.423/1.075,4.846/1.100}
      \draw[line width=0.6pt] (\x,-0.08) -- (\x,0)
        node[below=1.5pt] {$\lab$};
    \foreach \y/\lab in {0.000/1.04,1.283/1.075,2.567/1.110,3.850/1.145}
      \draw[line width=0.6pt] (-0.08,\y) -- (0,\y)
        node[left=1.5pt] {$\lab$};
    \node[below=9pt,font=\small] at (3.15,0) {$m_1$};
    \node[rotate=90,anchor=south,font=\small] at (-0.95,1.92) {$m_2$};

    % A=(1.072,1.066), B=(1.100,1.121)
    \coordinate (A) at (2.132,0.953);
    \coordinate (B) at (4.846,2.970);
    \draw[line width=1.45pt,black,-{Stealth[length=5pt]}] (A) -- (B);
    \fill[black] (A) circle (2.8pt);
    \draw[line width=0.95pt,black,fill=white]
      ($(B)+(-0.14,-0.14)$) rectangle ++(0.28,0.28);
    \node[anchor=east] at ($(A)+(-0.16,0.02)$) {A};
    \node[anchor=west] at ($(B)+(0.18,0.02)$) {B};
    \node[font=\scriptsize,align=left,anchor=north west] at (0.20,3.70)
      {continued mass segment};
  \end{scope}
\end{tikzpicture}
\caption{\label{fig:projection}Continuation connectivity and invariant projection.
(a)~Schematic of the two branches visible in the corrected scale-invariant
period--angular-momentum representation. A and B denote the documented pair
of distant mass points with nearly coincident invariant images; the panel is
not drawn to numerical scale. (b)~The straight segment in the $(m_1,m_2)$
plane followed by bidirectional continuation between the same endpoints,
whose mass separation is $0.062$. The connection shows that the projected
branch separation occurs within one sampled continuation component.}
\end{figure*}
